\section{Large girth outside a small set}\label{sec:half}

Recall that $q(G,P)$ is the probability that $\xi\cap E(G-P)$ is a forest.

\begin{theorem}\label{thm:half}
There are integers $D_*\ge7$ and $r_5$ with the following property. Let $G$ be a connected cubic
multigraph of rank $r\ge r_5$. Let $P$ be a nonempty set of at most $3\log r$ vertices such that
$G[P]$ is connected. If $G-P$ has girth at least $D_*$, then $q(G,P)<\frac25$.
\end{theorem}

\begin{idea}{outline of the proof}
Suppose that $q(G,P)\ge\frac25$. We find two disjoint parts of $G-P$ that are large, have small
boundaries and have large girth. Each of them has a forest proportion of at most $\frac35$, for
every choice of the edges outside it. This is a contradiction, because $(\frac35)^{2}<\frac25$. To
find the parts we use nested separators (Section~\ref{sec:separators}), a packing bound for
short cycles (Section~\ref{sec:packing}), and a second moment for cycles of moderate length
(Section~\ref{sec:halfproof}).
\end{idea}

\subsection{Pieces with a small boundary}\label{sec:affine}

Let $N$ be a connected subcubic multigraph. Its vertices of degree less than three are its
\emph{boundary vertices}. Let $T$ be a set of boundary vertices of even size. A \emph{$T$-join} of
$N$ is a set of edges in which the vertices of $T$ have odd degree and all other vertices have
even degree. There are $2^{\beta(N)}$ of them. The \emph{forest proportion} of $T$ is the
proportion of the $T$-joins that are forests. For $T=\emptyset$ the $T$-joins are the even
subgraphs, only the empty one is a forest, and the forest proportion is $2^{-\beta(N)}$. Put
\[
  R_{*}=R_{1/20},\qquad B=2R_{*}+2^{100}+1,
\]
with $R_{1/20}$ as in Section~\ref{sec:layers}. Then $\Theta_{1/20}(k)+k\le B(k+1)^{2}$ for all
$k\ge1$.

\begin{lemma}[Forest proportions at large girth]\label{lem:affine}
There are absolute constants $C_0$ and $C_1\ge1$ with the following property. Let $\Lambda\ge4$,
and let $N$ be a connected subcubic multigraph with at most $2\Lambda$ boundary vertices, at
least $C_1\Lambda^{4}$ vertices, and girth at least $C_0\log(\Lambda+2)$. Then every set of
boundary vertices of even size has forest proportion less than $\frac35$.
\end{lemma}

\begin{proof}
Let $s$ be the least integer with $2^{s}>(2s+1)B(2\Lambda+1)^{2}$. Then
$2^{s}\le2(2s+1)B(2\Lambda+1)^{2}$, and $B$ is an absolute constant. So there are absolute
constants $C_0$ and $C_1$ such that $2s<C_0\log(\Lambda+2)$ and $6\Lambda\,2^{s}<C_1\Lambda^{4}$
for every $\Lambda\ge4$.

Let $k$ be the number of boundary vertices of $N$. If $k\le1$, the only set is the empty one.
Then $N$ contains a cycle, and the forest proportion is at most $\half$. Let $k\ge2$, fix a
boundary vertex $o$, and call the set of the vertices at distance $j$ from $o$ the \emph{layer}
$j$.

\step{A large layer.} A ball of radius $s$ has at most $3\cdot2^{s}$ vertices. So at most
$3k\,2^{s}\le6\Lambda\,2^{s}$ vertices are within distance $s$ of a boundary vertex, and some
vertex $w$ has distance more than $s$ from every boundary vertex. All vertices within distance
$s$ of $w$ have degree three, and $2s$ is less than the girth. So the ball of radius $s$ around
$w$ is a tree in which every vertex at distance less than $s$ from $w$ has degree three, and it
has at least $2^{s}$ vertices. The distances from $o$ to two vertices of this ball differ by at
most $2s$. So the ball meets at most $2s+1$ layers, and some layer has at least
$2^{s}/(2s+1)>B(k+1)^{2}$ vertices.

\step{A marker.} Suppose that some set $T$ of boundary vertices has forest proportion at least
$\frac35$. Take a breadth-first tree of $N$ from $o$, and let $P$ be the union of its paths from
$o$ to the other boundary vertices. Then $N[P]$ is connected, every vertex outside $P$ has degree
three, and $P$ has at most $k$ vertices in each layer, because a path of the tree from $o$ has
one vertex in each layer that it meets. The tree paths inside $P$ contain a $T$-join $F_0$. The
map $J\mapsto J+F_0$ is a bijection from the $T$-joins of $N$ to the even subgraphs of $N$, and
$J+F_0$ has the same edges in $N-P$ as $J$. So if $J$ is a forest, then $J+F_0$ is a forest in
$N-P$. Hence $q(N,P)\ge\frac35>\half+\frac1{20}$. The distance from $o$ is a level function. By
Lemma~\ref{lem:width} with $\delta=\frac1{20}$, each layer has fewer than $\Theta_{1/20}(k)$
vertices outside $P$, and it has at most $k$ vertices in $P$. So each layer has fewer than
$B(k+1)^{2}$ vertices. This is a contradiction.
\end{proof}

\begin{lemma}[Product]\label{lem:product}
Let $G$ be a cubic multigraph, let $P$ be a set of vertices, and let
$U_1,\ldots,U_m\subseteq V(G)\sm P$ be pairwise disjoint sets such that each $G[U_i]$ is
connected. Suppose that in each $G[U_i]$ every even set of boundary vertices has forest
proportion at most $\frac35$. Then $q(G,P)\le(3/5)^m$.
\end{lemma}

\begin{proof}
Let $E'$ consist of the edges that lie in none of the graphs $G[U_i]$. Fix a restriction
$x\subseteq E'$ that extends to an even subgraph of $G$. For each $i$, let $T_i$ be the set of the
vertices of $U_i$ that are incident with an odd number of edges of $x$. The parity equations say
that $\xi\cap E(G[U_i])$ is a $T_i$-join. These equations involve disjoint sets of edges, so the
joins are independent and uniform, conditional on $\xi\cap E'=x$. A vertex of degree three in
$G[U_i]$ has no incident edge in $E'$. So $T_i$ is an even set of boundary vertices. If
$\xi\cap E(G-P)$ is a forest, its restriction to every $G[U_i]$ is a forest. The conditional
probability is therefore at most $(3/5)^m$. Average over $x$. Edges between the sets $U_i$ belong
to $E'$, so adjacent pieces cause no exception.
\end{proof}

From now on $D_{*}\ge7$ is a fixed integer such that $D\ge C_0\log(2D+2)$ for every $D\ge D_{*}$.

\subsection{Separators}\label{sec:separators}

\begin{idea}{two links give a separator}
If two disjoint connected sets have exactly two links, then one of them separates a neighbour of
the other one from a fixed vertex $p$. Disjoint sets that separate the same two vertices are
nested, and the parts between them have small boundaries (Figure~\ref{fig:chain}). So a long
chain of such sets contains two pieces to which Lemma~\ref{lem:affine} applies. We use this for
short cycles in Section~\ref{sec:packing} and for cycles of moderate length in
Section~\ref{sec:halfproof}.
\end{idea}

\begin{lemma}[Nested separators]\label{lem:nested}
Let $u$ and $p$ be two vertices of a connected graph $H$. Let $C_1$ and $C_2$ be disjoint connected
sets of vertices that avoid $u$ and $p$, such that $u$ and $p$ lie in different components of
$H-C_i$ for $i=1,2$. Let $U_i$ be the vertex set of the component of $u$ in $H-C_i$. Then exactly
one of the following holds: $C_1\subseteq U_2$ and $U_1\subseteq U_2$; or $C_2\subseteq U_1$ and
$U_2\subseteq U_1$.
\end{lemma}

\begin{proof}
The set $C_2$ is connected and disjoint from $C_1$, so it lies in one component of $H-C_1$.
Suppose that this component is not $U_1$. Then $U_1$ avoids $C_2$. Some edge joins $U_1$ to
$C_1$, because $H$ is connected. So $U_1\cup C_1$ is a connected set that contains $u$ and avoids
$C_2$, and therefore $U_1\cup C_1\subseteq U_2$. This is the first alternative. In the same way,
the second alternative holds if $C_1$ does not lie in $U_2$.

Suppose that $C_2\subseteq U_1$ and $C_1\subseteq U_2$. Let $W_i$ be the vertex set of the
component of $p$ in $H-C_i$. Then $W_2$ avoids $C_1$, and the argument above, with $p$ in the
place of $u$, gives $W_2\cup C_2\subseteq W_1$. This is impossible, because $C_2\subseteq U_1$,
and $U_1$ is disjoint from $W_1$. The two alternatives exclude each other, because
$C_1\subseteq U_2$ and $U_2\subseteq U_1$ would give $C_1\subseteq U_1$.
\end{proof}

So a family of pairwise disjoint connected sets that separate $u$ from $p$ is linearly ordered:
write $C_1<C_2$ in the first case of the lemma. The order is transitive, because it is the order
of the sets $U_i$ by inclusion.

\begin{figure}[ht]
\centering
\begin{tikzpicture}[x=1cm,y=1cm]
  % the cycle C and the vertex u
  \draw[hot] (0,0) ellipse [x radius=0.32, y radius=0.95];
  \node[note, text=warm] at (0,1.25) {$C$};
  \node[vtx] (u) at (0.85,0.25) {};
  \draw[edge] (0.31,0.25)--(u);
  \node[note, above=2pt] at (u) {$u$};
  % blocks
  \fill[accent!10] (2.4,-1.0) rectangle (5.2,1.0);
  \fill[accent!10] (5.8,-1.0) rectangle (8.6,1.0);
  \node[note, text=accent] at (3.8,-1.3) {a block};
  \node[note, text=accent] at (7.2,-1.3) {a block};
  % separating cycles
  \foreach \x in {2.1,5.5,8.9} \draw[cool] (\x,0) ellipse [x radius=0.26, y radius=0.95];
  \foreach \x in {2.9,3.5,4.1,4.7,6.3,6.9,7.5,8.1,9.7,10.3}
    \draw[accent, line width=0.45pt] (\x,0) ellipse [x radius=0.2, y radius=0.85];
  \node[note, text=accent] at (2.1,1.25) {$C_1$};
  \node[note, text=accent] at (5.5,1.25) {$C_{w+1}$};
  \node[note, text=accent] at (8.9,1.25) {$C_{2w+1}$};
  \node[note] at (10.95,0) {$\cdots$};
  % the marked set P
  \draw[draw=slate!70, fill=black!4, line width=0.6pt] (12.2,0) ellipse [x radius=0.55, y radius=0.45];
  \node[text=slate] at (12.2,0) {\small$P$};
  \node[vtx] (pp) at (11.92,0.05) {};
  \node[note, above=7pt] at (12.2,0.3) {$p\in P$};
  \draw[slate!60, line width=0.4pt, dash pattern=on 2pt off 2pt] (u)--(1.8,0.25) (11.3,0.05)--(pp);
\end{tikzpicture}
\caption{Nested separators. The vertex $u$ has a neighbour in a set $C$, and each set $C_i$
separates $u$ from $p$ in the graph without $C$. So the sets $C_i$ are nested
(Lemma~\ref{lem:nested}). Every $w$-th set is the end of a block, and the edges that leave a
block end in its two end sets (proof of Proposition~\ref{prop:chain}).}
\label{fig:chain}
\end{figure}

\begin{proposition}[A long chain]\label{prop:chain}
Let $G$ be a connected cubic multigraph, let $U$ be a set of vertices such that $G-U$ has girth at
least $D$, and let $P\subseteq U$ be a nonempty set such that $G[P]$ is connected. Let $p\in P$, let
$A$ be a region that avoids $P$, and let $u\notin A$ be a vertex in the component of $p$ in $G-A$.
Let $C_1,\dots,C_T$ be pairwise disjoint regions that avoid $P$, $A$ and $u$ and separate $u$ from
$p$ in $G-A$. Let $\Lambda\ge4$ be an integer and $m\ge1$ a number such that
$D\ge C_0\log(\Lambda+2)$, such that $A$ and every $C_i$ have at most $\Lambda$ boundary edges,
and such that every $C_i$ has at least $m$ vertices. If
\[
  T\;>\;\bigl(\Lambda+|U\sm P|+2\bigr)\bigl(\lceil2C_1\Lambda^{5}/m\rceil+1\bigr),
\]
then $q(G,U)\le\frac9{25}$.
\end{proposition}

\begin{proof}
Each $C_i$ is connected and separates $u$ from $p$ in $G-A$. So it lies in the component $H$ of
$G-A$ that contains $u$ and $p$. By Lemma~\ref{lem:nested} we may assume that
$C_1<C_2<\dots<C_T$. Let $U_i$ be the vertex set of the component of $u$ in $H-C_i$. Then
$U_i\cup C_i\subseteq U_j$ for $i<j$.

\step{Blocks.} For $i<j$ put $X_{ij}=U_j\sm(U_i\cup C_i)$. An edge of $H$ that leaves $U_j$ ends
in $C_j$, because $U_j$ is a component of $H-C_j$. No edge joins $X_{ij}$ to $U_i$: its end in
$X_{ij}$ would lie in the component $U_i$ of $H-C_i$. So every edge of $H$ that leaves $X_{ij}$
ends in $C_i$ or in $C_j$. If $i<l<j$, then $C_l\subseteq U_j$, and $C_l$ avoids
$U_i\cup C_i\subseteq U_l$; so $C_l\subseteq X_{ij}$. If $i<j\le i'<j'$, then
$X_{ij}\subseteq U_j\subseteq U_{i'}$ is disjoint from $X_{i'j'}$.

Put $w=\lceil2C_1\Lambda^{5}/m\rceil+1$, and let $Y_l=X_{lw+1,\,(l+1)w+1}$ for $l=0,1,\dots$ with
$(l+1)w\le T-1$. These sets are pairwise disjoint, each contains $w-1$ of the sets $C_i$, and
every edge of $H$ that leaves $Y_l$ ends in one of its two end sets. So
$|Y_l|\ge(w-1)m\ge2C_1\Lambda^{5}$, and at most $2\Lambda$ edges of $H$ leave $Y_l$.

There are at least $\Lambda+|U\sm P|+2$ sets $Y_l$. Discard $Y_l$ if it contains a vertex with a
neighbour in $A$, or a vertex of $U\sm P$. This discards at most $\Lambda+|U\sm P|$ sets, so two
sets remain. Let $Y$ be one of them. The edges of $G$ that leave $Y$ are edges of $H$, so
$|\delta Y|\le2\Lambda$. The set $P$ is connected, it contains $p$, and it is disjoint from $A$ and
from the sets $C_i$. So $P$ lies in the component of $p$ in $H-C_i$ for every $i$, and $Y$ avoids
$P$. Hence $Y$ avoids $U$, and $G[Y]$ has girth at least $D$.

\step{Pieces.} Every component of $G[Y]$ has an edge that leaves $Y$. So there are at most
$2\Lambda$ components, and one of them has a vertex set $Y'$ with at least $C_1\Lambda^{4}$
vertices. A boundary vertex of $G[Y']$ is an end of an edge of $\delta Y$, so there are at most
$2\Lambda$ of them. By Lemma~\ref{lem:affine}, every even set of boundary vertices of $G[Y']$ has
forest proportion less than $\frac35$. Lemma~\ref{lem:product}, for the two sets $Y'$ and the
marked set $U$, gives $q(G,U)\le\frac9{25}$.
\end{proof}

\begin{lemma}[Two links]\label{lem:sep}
Let $G$ be a connected multigraph, and let $p$ be a vertex. Let $\mathcal A$ be a finite family of
regions that avoid $p$, with weights $w_A\ge0$ of total $\mu>0$. Let $W_2$ be the sum of
$w_Aw_{A'}$ over the ordered pairs $(A,A')$ of disjoint regions of $\mathcal A$ with
$\ell(A,A')=2$. Let $N$ and $e$ be numbers such that, for every $A\in\mathcal A$, at most $N$
vertices outside $A$ have a neighbour in $A$, and the regions of $\mathcal A$ that are disjoint
from $A$ and joined to $A$ by an edge have total weight at most $e$. Let $K$ be a number with
$W_2\ge\frac23\mu^{2}-K\mu$ and $\mu>\frac32K+3e$. Then there are a region $A\in\mathcal A$ and a
vertex $u\notin A$ in the component of $p$ in $G-A$ such that the regions of $\mathcal A$ that
are disjoint from $A$, avoid $u$, and separate $u$ from $p$ in $G-A$ have total weight at least
$(\mu/3-K/2-e)/N$.
\end{lemma}

\begin{proof}
Let $A,A'\in\mathcal A$ be disjoint. Then $A'$ lies in the component of $p$ in $G-A$, or $A$ lies
in the component of $p$ in $G-A'$. Indeed, suppose that the first statement is false. A path from
$p$ to its first vertex in $A$ stays in the component of $p$ in $G-A$ before it ends, so it avoids
$A'$. As $A$ is connected and disjoint from $A'$, all of $A$ lies in the component of $p$ in
$G-A'$.

Call an ordered pair $(A,A')$ \emph{good} if the regions are disjoint, $\ell(A,A')=2$, no edge
joins them, and $A'$ lies in the component of $p$ in $G-A$. By the last paragraph, the pairs with
the first, the second and the fourth property have total weight at least $W_2/2$. The pairs that
are joined by an edge have total weight at most $e\mu$. So the good pairs have total weight at
least $W_2/2-e\mu\ge\mu\,(\mu/3-K/2-e)$, and for some $A\in\mathcal A$ the regions $A'$ with
$(A,A')$ good have total weight at least $\mu/3-K/2-e>0$. Fix this $A$.

Let $(A,A')$ be good. No edge joins the two regions, so the two links are components $Y_1,Y_2$ of
$G-A-A'$ that are joined to both. The vertex $p$ lies in at most one of them. Let $Y$ be one that
avoids $p$, and let $u$ be a vertex of $Y$ with a neighbour in $A$. The set $Y$ is connected and
joined to $A'$, and $A'$ lies in the component of $p$ in $G-A$. So $u$ lies in that component as
well. And $Y$ is a component of $G-A-A'$ that contains $u$ and not $p$. So $A'$ avoids $u$ and
separates $u$ from $p$ in $G-A$. There are at most $N$ vertices outside $A$ with a neighbour in
$A$. For one of them, $u$, the regions $A'$ have the total weight of the lemma.
\end{proof}

\subsection{Packing of short cycles}\label{sec:packing}

\begin{idea}{the purpose of the packing bound}
The second-moment argument requires a graph of girth about $\log r$. So we delete a maximal
family of disjoint short cycles, and we need a bound for the size of this family.
Lemma~\ref{lem:smallcuts} gives such a bound only when the proportion of forests exceeds one half
by a fixed margin. Here a large family has many pairs with exactly two links, as in the proof of
Lemma~\ref{lem:smallcuts}. By Lemma~\ref{lem:sep} they give a long chain of separators, and
Proposition~\ref{prop:chain} applies.
\end{idea}

\begin{proposition}[Packing]\label{prop:packing}
Let $G$ be a connected cubic multigraph, let $P$ be a nonempty set of vertices such that $G[P]$
is connected, and let $S$ consist of $s$ vertices outside $P$. Suppose that $G-(P\cup S)$ has
girth at least $D\ge D_*$ and that $q(G,P\cup S)\ge\frac25$. Let $d$ be an integer with
$4\le d\le2D$, and let $X_1,\ldots,X_n$ be pairwise disjoint cyclic regions outside $P\cup S$
with $|\delta X_i|\le d$ for all $i$. Then $n\le C_2\,d^{6}(s+1)\,2^{d}$, where
$C_2=44\,C_1$.
\end{proposition}

\begin{proof}
Put $U=P\cup S$.

\step{Many pairs with two links.} Let $E_i$ be the event that $\cO_{X_i}$ holds and
$\xi\cap E(G[X_i])\ne\emptyset$. Put $p_i=\PP(E_i)$, let $Z$ be the number of the events $E_i$
that occur, and put $\mu=\EE Z$. As in the proof of Lemma~\ref{lem:smallcuts}, we have
$p_i\ge2^{-d-1}$ and $\PP(E_i\cap E_j)=2^{\ell(X_i,X_j)-1}p_ip_j$ for $i\ne j$, and for each $i$
at most $d$ indices $j$ have $\ell(X_i,X_j)\ge3$. Let $W_2$ be the sum of $p_ip_j$ over the
ordered pairs $(i,j)$ with $\ell(X_i,X_j)=2$. Then $\EE Z^{2}\le\mu^{2}+W_2+(d+1)\mu$. If
$\xi\cap E(G-U)$ is a forest, then no event $E_i$ occurs. So $\PP(Z\ne0)\le\frac35$, and the
Cauchy--Schwarz inequality $\mu^{2}\le\EE Z^{2}\cdot\PP(Z\ne0)$ gives
\[
  W_2\;\ge\;\tfrac23\,\mu^{2}-(d+1)\,\mu .
\]
We have $n\le2^{d+1}\mu$. So the bound holds if $\mu\le6d$. Let $\mu>6d$.

\step{A chain.} Fix $p\in P$. We apply Lemma~\ref{lem:sep} to the sets $X_i$ with the weights
$p_i$. At most $d$ vertices outside $X_i$ have a neighbour in $X_i$, and at most $d$ of the other
sets are joined to $X_i$ by an edge. As $p_j\le1$, we may take $N=d$, $e=d$ and $K=d+1$. So
there are a set $X=X_i$, a vertex $u\notin X$ in the component of $p$ in $G-X$, and
$T\ge(\mu/3-2d)/d$ of the other sets that avoid $u$ and separate $u$ from $p$ in $G-X$. Each of
them has at most $d$ boundary edges. It contains a cycle of $G-U$, so it has at least
$D\ge d/2$ vertices. Also $D\ge C_0\log(d+2)$, because $d\le2D$ and $D\ge D_*$. As
$q(G,U)>\frac9{25}$, Proposition~\ref{prop:chain} with $A=X$, $\Lambda=d$ and $m=D$ gives
\[
  T\;\le\;(d+s+2)\bigl(\lceil2C_1d^{5}/D\rceil+1\bigr)\;\le\;(d+s+2)(4C_1d^{4}+2)
  \;\le\;7C_1\,d^{5}(s+1),
\]
because $d\ge4$. So $\mu\le3dT+6d\le22\,C_1d^{6}(s+1)$, and $n\le2^{d+1}\mu$.
\end{proof}

\paragraph{Raising the girth.}

\begin{theorem}[Raising the girth]\label{thm:raise}
Let $G$ be a connected cubic multigraph, and let $P$ be a nonempty set of vertices such that
$G[P]$ is connected. Suppose that $q(G,P)\ge\frac25$ and that $G-P$ has girth at least $D_*$. Put
$D_j=2^jD_*$ for $j\ge0$. There are nested sets $S_0\subseteq S_1\subseteq\cdots$ outside $P$
such that
\begin{enumerate}
\item[\textup{(i)}] $S_j$ is the union of a maximal family of pairwise disjoint cycles of length
at most $D_j$ in $G-P$; hence $G-(P\cup S_j)$ has girth greater than $D_j$;
\item[\textup{(ii)}] $\log(|S_j|+1)\le2D_j+C_3\log^{2}D_j$, where $C_3=7+\log(2C_2)$.
\end{enumerate}
\end{theorem}

\begin{proof}
Put $S_{-1}=\emptyset$.
Having chosen $S_{j-1}$, take a maximal disjoint family $\mathcal Q_j$ of cycles of length
at most $b=D_j$ in $G-(P\cup S_{j-1})$, and add its vertices to $S_{j-1}$ to form $S_j$.
The exterior girth is at least $a=D_*$ for $j=0$, and at least $a=D_{j-1}+1$ otherwise;
thus $b\le2a$.
Each cycle $C$ has $|\delta V(C)|\le|C|\le b$.
Also $q(G,P\cup S_{j-1})\ge q(G,P)\ge\frac25$.
Proposition~\ref{prop:packing} gives
\[
 |\mathcal Q_j|\le C_2\,b^{6}(|S_{j-1}|+1)2^b,
 \qquad |S_j|+1\le2C_2\,b^{7}2^b(|S_{j-1}|+1).
\]
The union of all families chosen so far is maximal at scale $D_j$, which proves~(i).
For~(ii), take logarithms in the recursion. The step $i$ adds at most
$D_i+7\log D_i+\log(2C_2)$. There are $j+1\le\log D_j$ steps, and $\sum_{i\le j}D_i<2D_j$. So
\[
 \log(|S_j|+1)\le2D_j+\log D_j\,\bigl(7\log D_j+\log(2C_2)\bigr)\le2D_j+C_3\log^{2}D_j .
\]
\end{proof}

\subsection{Proof of Theorem~\ref{thm:half}}\label{sec:halfproof}

Let $C_0$, $C_1$ and $D_*$ be as in Section~\ref{sec:affine}. Let $G$ and $P$ be as in the theorem, and
suppose that $q(G,P)\ge\frac25$. We derive a contradiction when $r$ is large. All limits below are
for $r\to\infty$, and they are uniform in $G$ and $P$. Put
\[
  x=\log r,\qquad A=6\,e^{124},\qquad L=\lfloor Ax\rfloor .
\]

\paragraph{Two girth scales.}
Let $D_j=2^{j}D_{*}$ and let $S_0\subseteq S_1\subseteq\cdots$ be the sets of
Theorem~\ref{thm:raise}. Let $D^{-}$ be the first $D_j$ that is at least $\sqrt x$, and let $D^{+}$
be the first $D_j$ that is at least $x/16$. For large $r$,
\[
  \sqrt x\le D^{-}<2\sqrt x,\qquad \frac x{16}\le D^{+}<\frac x8,\qquad D^{-}<D^{+}.
\]
Let $S^{-}\subseteq S^{+}$ be the corresponding sets, and put $J^{\pm}=G-(P\cup S^{\pm})$. Then
$J^{-}$ has girth greater than $D^{-}$, $J^{+}$ has girth greater than $D^{+}$, and
$J^{+}\subseteq J^{-}$. By Theorem~\ref{thm:raise}(ii),
\begin{equation}\label{eq:twosizes}
  |S^{-}|\;\le\;2^{4\sqrt x+C_3\log^{2}(2\sqrt x)}=r^{o(1)},\qquad
  |S^{+}|\;\le\;2^{x/4+C_3\log^{2}x}=r^{1/4+o(1)}.
\end{equation}
We use the small scale $D^{-}$ for the girth of the final pieces, and the large scale $D^{+}$
for a long chain.

\paragraph{Cycles at the large scale.}
We use Proposition~\ref{prop:loggirth} for $G$, $P$, $D=D^{+}$, $L$ and $S=S^{+}$. Its hypotheses
hold for large $r$: the set $P$ has at most $3x$ vertices and $G[P]$ is connected,
$x/16\le D^{+}\le x/8$, $7x\le L\le x\sqrt{\log x}$, and $|S^{+}|\le r/x^{2}$
by~\eqref{eq:twosizes}. So we are in the setting of Section~\ref{sec:moderate}, with $\Gamma=G$,
$U=P\cup S^{+}$ and $J=J^{+}$. The set $\mathcal C$ consists of the cycles of $J^{+}$ of length
at most $L$, and the numbers $s$, $a_C$, $\mu$ and $W_2$ and the random variable $Z$ are defined
there. With $\eta=r^{-1/60}$ the proposition gives $s\ge x/32-1$ and
\[
  L^{2}2^{-s}\le\eta,\qquad \mu\;\ge\;\frac14\ln\frac{L}{6x}-\frac1{\sqrt x}\;>\;30,\qquad
  \EE Z^{2}\;\le\;\mu^{2}+W_2+(1+\eta)\mu .
\]
The bound $\mu>30$ holds because $L/(6x)\ge e^{124}-1$.

\paragraph{Pairs with two links.}
If $\xi\cap E(G-P)$ is a forest, then $\xi$ contains no cycle of $J^{+}$, so $Z=0$. Hence $\PP(Z\ne0)\le\frac35$, and the
Cauchy--Schwarz inequality $\mu^{2}\le\EE Z^{2}\cdot\PP(Z\ne0)$ gives
\begin{equation}\label{eq:twolinks}
  W_2\;\ge\;\tfrac23\,\mu^{2}-(1+\eta)\,\mu .
\end{equation}

\paragraph{The chain.}
Fix $p\in P$. We apply Lemma~\ref{lem:sep} to the vertex sets of the cycles of $\mathcal C$, with
the weights $a_C$. A cycle $C\in\mathcal C$ has at most $L$ vertices, and each of them has one edge
end outside $C$. So at most $L$ vertices outside $C$ have a neighbour on $C$, and we may take
$N=L$. A cycle that is joined to $C$ by an edge passes through one of these vertices.
By~\eqref{eq:load} we may take $e=L^{2}2^{-s}\le\eta$, and $K=1+\eta$
by~\eqref{eq:twolinks}. As $\mu>30$ and $\eta\le\half$, the lemma gives a weight of more than
$8/L$. So there are a cycle $C\in\mathcal C$, a vertex $u\notin V(C)$ in the component of $p$ in
$G-V(C)$, and cycles of $\mathcal C$ of total weight more than $8/L$ that are disjoint from $C$,
avoid $u$, and separate $u$ from $p$ in $G-V(C)$. Take a maximal family $C_1,\dots,C_T$ of pairwise
disjoint cycles among them. Every cycle of the collection meets one of the at most $TL$ vertices
of the family. By~\eqref{eq:load}, $8/L\le TL\cdot L\,2^{-s}$. Hence
\[
  T\;\ge\;\frac{8\cdot2^{s}}{L^{3}}\;\ge\;\frac{4}{A^{3}}\cdot\frac{r^{1/32}}{x^{3}}.
\]

\paragraph{The contradiction.}
The vertex set of a cycle of $J^{+}$ of length at most $L$ has at most $L$ boundary edges and more
than $D^{+}\ge x/16$ vertices. We apply Proposition~\ref{prop:chain} with $U=P\cup S^{-}$,
$D=D^{-}$, $A=V(C)$, $\Lambda=L$ and $m=x/16$. The graph $G-U=J^{-}$ has girth greater than
$D^{-}$, and $D^{-}\ge\sqrt x\ge C_0\log(L+2)$ for large $r$. By~\eqref{eq:twosizes},
\[
  \bigl(L+|S^{-}|+2\bigr)\bigl(\lceil32C_1L^{5}/x\rceil+1\bigr)\;=\;r^{o(1)}\;<\;T
\]
for large $r$. So $q(G,P)\le q(G,P\cup S^{-})\le\frac9{25}<\frac25$. This contradiction proves
Theorem~\ref{thm:half}. \qed
